\BourbakiBackSection{BIBLIOGRAPHY}

{[}1{]} O. NEUGEBAUER, Vorlesungen über Geschichte der antiken Mathematik, Bd. I: Vorgriechische Mathematik, Berlin (Springer), 1934.

{[}2{]} Euclidis Elementa, 5 volumes, ed.~J. L. Heiberg, Lipsiae (Teubner), 1883-88.

{[}2a{]} T. L. HEATH, The Thirteen Books of Euclid's Elements\ldots, 3 volumes, Cambridge, 1908.

{[}3{]} Archimedis Opera Omnia, 3 volumes, ed.~J. L. Heiberg, 2nd edition, 1913-15.

{[}3a{]} Les Œuvres complètes d'Archimède, translated by P. Ver Eecke, Paris-Bruxelles (Desclée-de Brouwer), 1921.

{[}4{]} Diophanti Alexandrini Opera Omnia\ldots, 2 volumes, ed.~P. Tannery, Lipsiae (Teubner), 1893-95.

{[}4a{]} Diophante d'Alexandrie, translated by P. Ver Eecke, Bruges (Desclée de Brouwer), 1926.

{[}5{]} R. BOMBELLI, L'Algebra, ed.~E. Bortolotti, Bologna (Zanichelli), 1929.

{[}6{]} LES ŒUVRES Mathématiques de SIMON STEVIN de Bruges, Ou sont insérées les MÉMOIRES MATHÉMATIQUES, Esquelles s'est exercé le Très-haut et Très-illustre Prince MAURICE DE NASSAU, Prince d'Aurenge, Gouverneur des Provinces des Païs-Bas unis, General par Mer et par Terre, etc., Le tout reveu, corrigé et augmenté par ALBERT GIRARD Samielois, Mathématicien, A LEYDE, Chez Bonaventure et Abraham Elsevier, Imprimeurs ordinaires de l'Université, Anno MDCXXXIV (= 1634), vol.~I.

{[}7{]} I. BARROW, Mathematical Works, Cambridge (University Press), 1860.

{[}8{]} I. NEWTON, De Analysi per equatione numero terminorum infinitas, in Commercium Epistolicum D. Johannis Collins et aliorum de Analysi promota, Londini, 1712.

{[}9{]} C. F. GAUSS, Werke, vols. III (Göttingen, 1816) and \(\mathbf{X}_1\) (ibid., 1917).

{[}10{]} A. CAUCHY, Cours d'Analyse de l'École Royale Polytechnique, 1re partie, 1821 = Œuvres (II) vol.~III, Paris (Gauthier-Villars), 1897.

{[}11{]} B. BOLZANO, Rein Analytischer Beweis des Lehrsatzes, dass zwischen je zwei Werthen, die ein entgegengesetztes Resultat gewähren, wenigstens eine reelle Wurzel liegt, Ostwald's Klassiker, no. 153, Leipzig, 1905.

{[}12{]} N. H. ABEL, Œuvres, 2 volumes, ed.~Sylow and Lie, Christiania, 1881.

{[}13{]} R. DEDEKIND, Gesammelte mathematische Werke, vol.~II, Braunschweig (Vieweg), 1932, p.~315.

{[}14{]} E. BOREL, Leçons sur la théorie des fonctions, 2nd edition, Paris (Gauthier-Villars), 1914.

{[}15{]} R. BAIRE, Leçons sur les fonctions discontinues, Paris (Gauthier-Villars), 1905.

{[}16{]} N. LUSIN, Leçons sur les ensembles analytiques et leurs applications, Paris (Gauthier-Villars), 1930.

